\section{Discussion}
\label{sec:discussion}

\subsection{A middle position, and why it is a position rather than a hedge}

Two literatures reach opposite conclusions about what intergenerational
transmission of fertility implies. \citet{collins2019} apply the breeder's
equation to national fertility and conclude that global population stabilisation
this century becomes very unlikely; the European and North American trajectories
in their model turn from decline to growth. \citet{arenberg2022} reply that a
model built from transmission alone contains no absolute fertility level, cannot
represent the fact that higher-fertility subgroups are themselves converging
downward, and therefore cannot establish that transmission produces growth at
all.

Both are right about what they are arguing. The disagreement is not really about
the mechanism, which nobody disputes; it is about whether a model that carries
only a transmission coefficient can support a conclusion about population size.
It cannot, and the fix is to put the mechanism inside a projection that carries
levels.

Doing that yields a result that neither side would predict from the other's
setup. The mechanism is real and measurable, its direction is unambiguous, and it
is worth \LadderMainstreamGain{} billion people at \ProjectionEndYear{} --- far
too large to leave out of a long-run projection. It is also not a return to
growth. Under the benchmark environment with mainstream selection, the world
still peaks and then runs close to flat; what selection changes is where the
curve settles. The Arenberg condition is satisfied rather than violated, because
absolute age-specific fertility is carried at every step, and the outcome is a
level shift of about \LadderMainstreamGainPercent{}.

The second finding is more surprising and cuts against the way this subject is
usually discussed. Explicitly named high-fertility groups contribute
\NamedGroupShareOfMechanism{} of what unlabelled variation contributes, even
though the model is generous to them at every choice point. The compositional
force in long-run fertility is not a story about visible minorities; it is a
property of the ordinary dispersion of family size, acting on everybody at once.
That is a better result than the alternative, because ordinary dispersion is a
quantity that national statistical offices already measure.

\subsection{On the environment, and why we report a threshold}

The honest position on the fertility environment is that nobody knows what it
will do. The literature contains coherent arguments in both directions: the
low-fertility trap hypothesis \citep{lutz2006} describes self-reinforcing
downward pressure through ideal family size and cohort experience, while the
mean-reverting specifications used in official projections
\citep{alkema2011,raftery2012} embed an eventual recovery. Neither is settled,
and the sensitivity of a 2150 world total to the difference is enormous ---
\LadderPressureLoss{} billion for the illustrative path used here.

A projection whose headline is controlled by that quantity is reporting its
author's intuition with extra steps. \Cref{eq:break-even} is the alternative:
because both forces multiply fertility, the model can be asked for the
environmental decline that exactly offsets the measured compositional gain, and
that number --- \BoundaryPaired{} per decade after \PressureFromYear{}, robust to
migration to within \BoundaryPairedLow{}--\BoundaryPairedHigh{} --- is determined
by the two measured parameters and the arithmetic.

We would defend this as the right output for a problem of this shape. It is
falsifiable in the parameters that are measurable, it is explicit about the one
that is not, and it turns a disagreement about the future into a disagreement
about a rate. Someone who believes the fertility environment will keep
deteriorating rapidly and someone who believes it has nearly bottomed out can
now locate their disagreement on \Cref{fig:boundary} rather than trading
projections.

\subsection{Limitations}
\label{sec:limitations}

\begin{description}[leftmargin=1.4em,style=nextline]

\item[The environment multiplies every type equally.] This is what makes the
boundary a single clean number, and it is the assumption most likely to be
wrong. Rising costs of housing, childcare and foregone earnings plausibly fall
hardest on exactly the people whose disposition would otherwise produce large
families, in which case the environment would interact with selection rather
than merely oppose it and no single threshold would exist. Modelling that
interaction requires evidence on how fertility declines have differed across the
family-size distribution within cohorts, which we did not find at the coverage
needed. We regard this as the most valuable single extension.

\item[Three types are a discretisation, not a discovery.] The propensity
distribution is stood in for by \MainstreamTypes{} points chosen to reproduce its
first two moments exactly. Higher moments, and the shape of the upper tail, are
structural approximations. The number of types is a declared scenario knob for
that reason.

\item[The dispersion parameter is an effective spread, not a latent trait.] The
concession to \citet{hruschka2016} in \Cref{sec:calibration} is genuine. What is
defended is the moment match to an observed covariance; what is not claimed is
that the observed variance in completed family size is inherited. Because the
covariance is what enters \Cref{eq:response}, this is sufficient for the
one-generation response, and it is the extrapolation across four generations that
the empirical range \FamilySizeCVLow{}--\FamilySizeCVHigh{} is carried to reflect.

\item[Transmission is one constant, applied everywhere.] A single parent--child
correlation is used for all 237 countries, when the underlying literature
reports regional means ranging from 0.09 to 0.19 and does not cover most of the
world. Countries outside Europe and North America are effectively assigned a
parameter estimated inside them.

\item[Named groups are modelled without conversion in, and with their host's age
structure.] Both push the same way, so the reported group shares are lower
bounds, but neither is small in principle and a group large enough to matter
would deserve its own age structure.

\item[After 2100 the demographic schedules are held constant.] This is our rule
and not the source's, and \Cref{sec:horizon} measures what it is worth
(\WidthHoldConstant{} billion of world spread at \ProjectionEndYear{}) rather
than leaving it implicit.

\item[Migration by age is borrowed.] The United Nations does not publish net
migrants by single year of age. Where migration is an input it is a residual
backed out of the UN's own path, which is usable but is not independent evidence;
where migration is the object of study its level comes from published
trajectories but its age and sex composition is still borrowed.

\item[Selection is on fertility only.] Survival is identical across types.
High-fertility groups often have distinctive mortality, and modelling that
without evidence would add a parameter and subtract credibility.

\end{description}

\subsection{What would show this to be wrong}
\label{sec:falsification}

The mechanism's claim resolves in cohort fertility, not period fertility, and the
distinction is not pedantry: period rates are depressed by postponement even when
completed family sizes are unchanged \citep{bongaarts1998}, so a mechanism
predicting a rise in completed cohort fertility could be scored as failing for
decades by a measure that is not measuring it.

The implication is uncomfortable but should be stated: the first genuinely
informative test of this model arrives when the cohorts born around 1990 complete
their childbearing, in the late 2030s. What can be checked before then is
narrower --- whether the dispersion of completed family size behaves as
\Cref{sec:cv} measures it in countries not yet in the sample, whether the
parent--child correlation holds up outside Europe and North America, and whether
the composition drift in \Cref{fig:mechanism}b is visible in successive cohorts
of the same country.

Because a claim that resolves in 2038 is easy to revise quietly in 2035, the
projections behind this paper are written as immutable dated records with proper
scoring rules attached \citep{gneiting2007}, and the storage layer raises rather
than overwrite an existing record. The comparator of \Cref{sec:comparator} was
deposited before the mechanism results existed, and is marked as not a claim of
this project. We think this record-keeping, rather than any particular number in
\Cref{sec:results}, is the part of the work most likely to still be useful in
thirty years.
